% Substitua as instrucoes pelo resumo do portfolio concluido.
\undfsetcontent{resumo}{%
Este texto orienta a elaboração do resumo do relatório final de um portfólio
de Projetos Aplicados. Apresente o percurso acadêmico abrangido, seu contexto
e a questão que articula o conjunto de projetos. Indique o objetivo do
portfólio e explique como os registros parciais e as evidências foram
recuperados, organizados e analisados. Identifique os tipos de projetos
realizados e a contribuição individual, distinguindo as atividades em equipe.
Sintetize as principais entregas e os resultados efetivamente sustentados
pelos materiais disponíveis. Apresente as relações entre os projetos, as
decisões tomadas, as dificuldades e os aprendizados discutidos na análise
integradora. Explicite limites de memória, de documentação e de avaliação,
sem transformar impressões ou expectativas em resultados demonstrados. Conclua
com a contribuição do conjunto para a trajetória analisada. O resumo deve
representar o relatório completo, e não descrever somente um projeto ou
reproduzir a sequência dos registros parciais. Redija um parágrafo único,
sem citações bibliográficas, e substitua integralmente estas orientações
antes da entrega. Confira a correspondência com a versão em língua estrangeira
e escolha palavras-chave que representem o conteúdo final do portfólio.}
\undfsetcontent{palavraschave}{portfólio; projetos aplicados; análise integradora}
\undfsetcontent{abstract}{%
This paragraph explains how to summarize a completed Applied Projects portfolio.
Present the academic trajectory covered by the report, its context and the
question connecting the projects. State the portfolio objective and explain
how partial records and supporting evidence were recovered, organized and
examined. Identify the kinds of projects undertaken and the student's own
contribution, distinguishing individual work from team activities. Summarize
the main deliverables and the findings actually supported by the available
materials. Describe connections between projects, decisions, difficulties and
learning discussed in the integrative analysis. Make limitations of memory,
documentation and evaluation explicit, without presenting expectations or
impressions as demonstrated outcomes. Conclude with the contribution of the
complete set of projects to the trajectory examined. The abstract should
represent the entire final report rather than describe a single project or
repeat a sequence of partial records. Use one coherent paragraph without
bibliographic citations and replace these instructions before submission.
Check that the Portuguese and foreign language versions describe the same
scope, procedures, findings and conclusions, and choose keywords that accurately
represent the completed portfolio.}
\undfsetcontent{keywords}{portfolio; applied projects; integrative analysis}
